\section{Total asphericity and ordinary colimits}
\begin{claim}{F.1}{Proved}{The action criterion used in this part}
Here total asphericity means the condition of XII ms.~88 tested by all ordinary final images; we do not substitute the different diagonal-asphericity terminology used elsewhere in the literature. For a nonempty connected small category $A$, total $W_\infty$-asphericity is equivalent to every connected component of $\operatorname{El}_A(H)$ being weakly contractible, for every covariant $H:A\to\mathrm{Set}$. It is also equivalent to the same condition with integral acyclicity in place of weak contractibility. Exactness of $\colim_A:\Ab^A\to\Ab$ implies these conditions.
\end{claim}
\begin{proof}
If $A$ is totally integrally aspherical, it is integrally acyclic by the identity test. For a connected nonempty category of elements $E=\operatorname{El}_A(H)$, adjoin an object $v$ to $A$ with arrows $v\to a$ the elements of $H(a)$, and composition induced by $H$. The inclusion $A\to A_H$ is final: its comma at $v$ is $E$, and the other commas have initial objects. Thus $BA_H$ is acyclic. Its nerve is the homotopy pushout $BA\leftarrow BE\to *$. The reduced homology sequence gives acyclicity of $BE$. For a component of an arbitrary action, restrict $H$ to that component.

This property passes to every component $E$ of an action category. Indeed for $K:E\to\mathrm{Set}$ the action $L(a)=\coprod_{h\in H_E(a)}K(a,h)$ on $A$ has category of elements exactly $\operatorname{El}_E(K)$, with all its morphisms. Now take $G=\pi_1BE$. For a cyclic subgroup $C\leq G$, the action on cosets in the groupoid completion of $E$ defines a connected covering category with fundamental group $C$. Its first integral homology must vanish by the inherited property, but that homology is $C$. All cyclic subgroups are therefore trivial. Hence $BE$ is simply connected and acyclic, and homological Whitehead makes it weakly contractible.

Conversely every comma $b\backslash A$ of an ordinary final functor $A\to B$ is a connected nonempty action category, for the action $a\mapsto B(b,f(a))$. The assumption and Quillen's Theorem A make $BA\to BB$ a weak equivalence. The singleton action makes $BA$ contractible. Thus every final image is contractible, as required.
Finally, for any $F:\operatorname{El}_A(H)\to\Ab$, the functor
$(\Sigma_HF)(a)=\bigoplus_{h\in H(a)}F(a,h)$ is exact and has
$\colim_A\Sigma_HF=\colim_{\operatorname{El}_A(H)}F$.
Exactness on $A$ implies exactness on its action categories and on their components. Apply this to the usual representable bar resolution of the constant integral module; its derived colimits are integral nerve homology. This proves the last assertion, and also specifies the variance of every action used.
\end{proof}

\begin{claim}{F.2}{Proved}{What totality forces before the affine gap}
A totally $W_\infty$-aspherical category has common successors for every finite family of objects. For every original parallel pair $a,b:p\to q$, there are an object $r$ and $t\in\mathbb Z[A(q,r)]$ of augmentation one such that $ta=tb$.
\end{claim}
\begin{proof}
If two objects have no common successor, their upward closures $U,V$ are disjoint cosieves. Start with the constant two-point action and identify the two points exactly on $U\cup V$. Its element category is the double of $A$ glued along $U\amalg V$. Nerves form the corresponding homotopy pushout. Since $BA$ is contractible and both $U,V$ are connected, the Mayer--Vietoris sequence gives a nonzero first homology group of the double, contrary to F.1. Induction gives finite common successors.
For the parallel pair form the quotient set action $H$ of $A(p,-)$ generated by $a\sim b$. In the exact sequence
\[
0\longrightarrow I\longrightarrow\mathbb Z[A(p,-)]\longrightarrow\mathbb Z[H]\longrightarrow0,
\]
$I$ is generated by $\delta=a-b$ at $q$. The middle diagram is projective, and F.1 gives $L_1\colim\mathbb Z[H]=0$. The image of $\colim I$ in $\colim\mathbb Z[A(p,-)]=\mathbb Z$ is zero, so $\colim I=0$. Write the element $\delta$ as a finite sum of the relations $I(v_j)x_j-x_j$ in the direct-sum presentation of this colimit. Choose a common successor $r$ for their finitely many objects and arrows $d_s:s\to r$. Write $x_j=\sum_l n_{jl}c_{jl}\delta$, with $c_{jl}:q\to\operatorname{dom}(v_j)$. Then
\[
t=d_q-\sum_{j,l}n_{jl}(d_{\operatorname{cod}(v_j)}v_j-d_{\operatorname{dom}(v_j)})c_{jl}
\]
has augmentation one and satisfies $t\delta=0$. All terms are arrows $q\to r$, so this identity is in the specified free abelian group.
\end{proof}

\begin{claim}{F.3}{Proved}{The requested injection and free-category examples}
Let $A_S$ be the full subcategory of $\Delta_{\mathrm{inj}}$ on $[n]$, $n\in S\subseteq\mathbb N$. It is totally $W_\infty$-aspherical exactly when $S$ is a singleton or unbounded. In these cases its integral colimits are exact. The category of finite sets and injections, with or without the empty set, is not totally $W_\infty$-aspherical. The one-object category of injective self-maps of $\mathbb N$ is not totally $W_\infty$-aspherical. For the free category on any small directed graph, total $W_\infty$-asphericity is equivalent to filteredness; local finiteness is unnecessary.
\end{claim}
\begin{proof}
We first give the affine cones needed for unbounded $S$. For each $n\geq0$ and $k\geq1$, there is a homogeneous integral polynomial $C_{n,k}(z_0,\ldots,z_{n+1})$ of value one at $(1,\ldots,1)$ such that $y_i^kD_iC_{n,k}$ is independent of $i$, where $D_i$ identifies $z_i,z_{i+1}$ and relabels the variables as $y_0,\ldots,y_n$. For $n=0$ take $1$. Inductively set $Q=C_{n-1,2k}$ and prescribe
$D_iC_{n,k}=Q\prod_{l\ne i}y_l^k$.
The prescriptions agree on pairwise intersections: for adjacent merged pairs this is immediate; for separated pairs the equality is precisely $w_q^{2k}D_qQ=w_p^{2k}D_pQ$ from the induction hypothesis, multiplied by the common remaining factors. In coordinates $x_i=z_i-z_{i+1}$ and $u=z_{n+1}$, these are compatible restrictions to the coordinate hyperplanes $x_i=0$. If $\Phi_T$ is their common restriction to $x_i=0$, $i\in T$, lift it to the full polynomial ring by making it independent of the omitted coordinates $x_i$. Then inclusion--exclusion
$C=\sum_{\varnothing\ne T}(-1)^{|T|+1}\Phi_T$
extends them. Substitution back into $z$ preserves integrality and homogeneity. Its degree is $k(2^{n+1}-n-2)$.

An injection $g:[n]\to[M]$ corresponds to its gap monomial
\[
 z_0^{g(0)}z_1^{g(1)-g(0)-1}\cdots
 z_n^{g(n)-g(n-1)-1}z_{n+1}^{M-g(n)}.
\]
Deleting $g(i)$ corresponds to $P\mapsto y_iD_iP$. Thus $C_{n,1}$ gives an integral affine arrow $t:[n]\to[2^{n+1}-2]$ for $n\geq1$ equalizing all cofaces. More explicitly, after $r$ deletions the common gap polynomial is
\[
 \Bigl(\prod_j z_j^{\,2^r-1}\Bigr)C_{n-r,2^r}.
\]
One further deletion identifies two exponents and contributes one additional gap, so its varying factor is $y_i^{2^r}D_iC_{n-r,2^r}$, independent of $i$ by construction. The formula therefore follows by induction, for every deletion sequence. Every injection into $[n]$ is such a sequence of cofaces, so $t$ equalizes every parallel pair. If $S$ is unbounded, postcompose with an injection into some $[N]$, $N\in S$, above this target. Hence the category $\operatorname{Aff}_{\mathbb Z}(A_S)$ is filtered: it has common successors and the same $t$ equalizes every parallel affine pair, because both have total coefficient one. Extend any abelian diagram linearly to this affine category. Its cocones are unchanged, and filtered colimits are exact. F.1 proves totality. The polynomial identities and all parallel injections are checked independently for $n\leq3$ by the program; the induction proves all $n$.
A singleton $S$ gives the terminal category. If $S$ is finite with at least two elements and maximum $N$, there are distinct injections $[m]\to[N]$ from some smaller $m\in S$, whereas the only arrow out of $[N]$ is its identity. They cannot satisfy F.2. The empty category is not total. This proves the complete stated classification.

For finite sets let $H(X)$ be the set of two-element subsets of $X$. The category of elements consists of $(X,T)$, $|T|=2$, with injections carrying $T$ onto the target distinguished subset. The functor $(X,T)\mapsto(T,T)$ retracts it onto the groupoid of two-element sets. The inclusions $T\hookrightarrow X$ give a natural transformation from the inclusion of this retract to the identity. Thus the entire element category has type $BC_2$, and is connected with $H_1=\mathbb Z/2$. This contradicts F.1. The finite restriction through cardinality three is separately checked by its complete edge-path presentation and Tietze eliminations.

For the injection monoid, let $s(n)=n+1$. An affine equalizer of $1$ and $s$ would be an invariant finite-support integral signed measure of mass one for $f\mapsto f s$ on injective self-maps. Such a self-map has no periodic point, since $f(n+r)=f(n)$ is impossible for $r>0$. An invariant finite-support signed measure is supported on finite cycles: each point in its finite support must have a predecessor in that support, and counting predecessors forces the self-map to permute the support. No such measure exists. F.2 excludes totality.

In a free category, no distinct parallel paths admit an affine equalizer. If their lengths differ, the terms of maximal length in the product with the longer path cannot cancel. If their lengths agree, concatenation with the two different initial words has disjoint images in the path basis. Thus F.2 forces the free category to have at most one arrow between any two objects. Its common-successor property then makes it filtered. Conversely filtered categories have exact integral colimits and hence are total by F.1. This proves the last assertion in its stated generality.
\end{proof}

\begin{claim}{F.4}{Proved}{Weakly terminal objects and the finite-object case}
A totally $W_\infty$-aspherical category with a weakly terminal object is filtered. Consequently the universal Cat-colimit property in Part F(3) implies filteredness for every category with a weakly terminal object, and for every nonempty connected category with finitely many objects.
\end{claim}
\begin{proof}
We include the monoid step because signed coefficients cannot be replaced by positive ones. Suppose a monoid $M$ has, for every $x,a$, an affine $t$ with $tx=txa$. The measure $tx$, of mass one and supported in $Mx$, is invariant under right multiplication by $a$. Its support is a union of finite periodic orbits and its coefficients are constant on each orbit, by the support-counting argument in F.3. Choose the least attainable positive period $r$ in $Mx$. If a prime $\ell$ divides $r>1$, choose $z\in Mx$ of period $r$ and apply the hypothesis to $z,a$. The resulting invariant measure on $Mz$ has mass one, so one of its orbit lengths $s$ is not divisible by $\ell$. But an element of that orbit is also fixed by $a^r$, since it lies in $Mz$. Its period divides $\gcd(r,s)<r$, a contradiction. Thus some $y$ satisfies $yxa=yx$. Apply this with $x=1$ and then with $x=y$ to equalize any given pair of elements of $M$. Hence $M$ is filtered.

If $q$ is weakly terminal in $A$, postcompose the affine equalizers supplied by F.2 back to $q$. The preceding lemma makes $\operatorname{End}_A(q)$ filtered. For $a,b:p\to r$, choose $d:r\to q$ and an affine endomorphism $t=\sum n_ic_i$ of $q$ equalizing $da,db$. Filteredness of the endomorphism monoid gives $z$ with $zc_i=z$ for every $i$. Since $\sum n_i=1$, the genuine arrow $zd:r\to q$ equalizes $a,b$. Together with common successors this proves filteredness.

Now let $I$ be nonempty and connected, and assume every diagram $I\to\Cat$ has its ordinary colimit computed by the canonical Grothendieck-construction comparison. Apply this to discrete-valued diagrams $H:I\to\mathrm{Set}$. It says that every component of $\operatorname{El}_I(H)$ is contractible. The constant singleton diagram also makes $BI$ contractible. F.1 proves totality, so the preceding result applies when $I$ has a weakly terminal object. In the finite-object case common successors provide such an object. Conversely filtered shapes satisfy the universal comparison: every finite list of objects, arrows, and relations in an ordinary category colimit is represented at a common stage, and every comparison comma category is filtered. Quillen's Theorem A proves the specified comparison is a weak equivalence.
\end{proof}

\begin{claim}{F.5}{Open}{The exact connected converse}
For a nonempty connected small category $A$ with no weakly terminal object, we do not settle whether total $W_\infty$-asphericity forces exactness of $\colim_A:\Ab^A\to\Ab$. The exact gap is whether the hypotheses of F.2 equalize every parallel pair in $\operatorname{Aff}_{\mathbb Z}(A)$, not merely every original parallel pair. F.3 resolves the listed injection and free-category candidates but does not remove this gap for a general countable category. Exact integral colimits remain a proved sufficient condition, and the affine-filtered criterion is the precise algebraic reformulation. The original claims of equality with filteredness are false by the unbounded $A_S$ examples.
\end{claim}

\begin{claim}{F.6}{Open}{Total presheaves and composition of dwyerienne sieves}
We do not settle whether every totally $W_{\mathbb Z}$-aspherical presheaf is totally $W_\infty$-aspherical, or equivalently the composition assertion specified in TOTASPH F.27. The proof of F.1 does not transfer: for a presheaf $H$ on $A$, totality tests only those covariant actions on its element category pulled back from $A$, whereas the cyclic-cover argument needs arbitrary actions on that element category. The search over all monoids of order at most six and all presheaves of cardinality at most six was not performed in this return. The smaller monoid catalogue in Part A is not such a search and cannot certify this assertion. No acyclic-group example satisfying the required universal restricted-action test is provided here.
\end{claim}

\begin{claim}{F.7}{Open}{Universal Cat-colimit shapes without a weakly terminal object}
For a connected countable category $I$ without a weakly terminal object, and hence also in general, the implication from the universal ordinary-Cat-colimit comparison to filteredness remains unresolved here. F.4 gives a necessary totality condition and the full weakly-terminal case. The unbounded injective-ordinal examples do not refute this implication: the diagram $[n]\mapsto B\mathbb N^{n+1}$ has categorical colimit $B\mathbb N\simeq S^1$, whereas its Grothendieck construction has integral homology $\mathbb Z$ in every degree. On all of $\Delta_{\rm inj}$ its degree-$k$ fibre homology, for $k\geq1$, is the projective representable $\mathbb Z[\Delta_{\rm inj}([k-1],-)]$; the homotopy-colimit spectral sequence collapses. On an arbitrary unbounded full subcategory $A_S$, the restricted diagram need not be representable. Instead its higher derived colimits vanish by exactness in F.3, and its colimit is $\mathbb Z$: any two marked $k$-element ordered subsets admit embeddings into a common larger ordinal identifying their marks, then into an object of $A_S$. The categorical colimit is still $B\mathbb N$, since the same embeddings identify all coordinate generators and introduce precisely their commutative-monoid relation. Thus all these shapes fail the premise. This prevents importing the nonfiltered total example as an answer to the different Cat-colimit question.
\end{claim}
